\subsection{迷向子空间。}

定义 1. --- 给定环 A 上的模 E，若 \(\Phi(x, x) = 0\)，则称 E 的元素 x 是迷向的。称 E 的子模 F 是

\begin{enumerate}
\def\labelenumi{\arabic{enumi})}
\item
  迷向的，如果存在 F 中与 F 正交的元素 \(x \neq 0\)；
\item
  全迷向的，如果 \(\Phi\) 在 F 上的限制为零。
\end{enumerate}

当模 E 装备了二次型 Q 时，若 E 的元素关于与 Q 相伴的双线性型是迷向的，则称该元素是迷向的（相应地，若 E 的子模关于与 Q 相伴的双线性型是迷向的或全迷向的，则称该子模是迷向的或全迷向的）。

迷向向量就是与自身正交的向量。说子模 F 是迷向的，就意味着 \(F \cap F^{0} \neq \{0\}\)，或者等价地，\(\Phi\) 在 F 上的限制是退化的；因此 E 的非迷向子模 G 是使 \(\Phi\) 在 G 上的限制非退化的子模。要使 E 的子模 F 是全迷向的，必须且只需 \(F \subset F^{0}\)。若 F 是 E 的全迷向子模，则包含在 F 中的每个子模 \(F'\) 也是全迷向的。两两正交的全迷向子模族之和是全迷向子模。E 的全迷向子模的集合按包含关系排序显然是归纳的；由此推出，每个全迷向子模都包含在某个极大全迷向子模之中。

命题 1. --- 假设 A 是一个体。设 F 是 E 的有限维非迷向子空间；则 E 是 F 与 F⁰ 的直和。

事实上，由于按假设 \(\Phi'\) （\(\Phi\) 在 F 上的限制）非退化，故映射 \(d_{\Phi'}\) （F 到其对偶 F* 的、右相伴于 \(\Phi'\) 的）是单射，又因 F 与 F* 是相同有限维的两个空间，故它是双射。于是，对每个 \(y \in E\)，存在 F 中唯一的元素 \(y_0\)，使得 \(\Phi(x, y) = \Phi(x, y_0)\) 对所有 \(x \in F\) 成立，亦即 \(y - y_0 \in F^0\)；这就证明了 E 是 F 与 F \(^0\) 的直和。

推论. --- 若 F 是 E 的有限维子空间，且 Φ 非退化，则下列条件等价：

\begin{enumerate}
\def\labelenumi{\alph{enumi})}
\item
  F 非迷向。
\item
  \(\mathbf{F}^0\) 非迷向。
\item
  E 是 F 与 F \(^{0}\) 的直和。
\end{enumerate}

命题 1 确实表明 \(a)\) 蕴含 \(c)\)，而 \(c)\) 蕴含 \(a)\) 与 \(b)\)。最后，若 \(\mathbf{F}^0\) 非迷向，则有 \(\mathbf{F} \cap \mathbf{F}^0 \subset \mathbf{F}^0 \cap \mathbf{F}^{00} = \{0\}\)；于是 \(\mathbf{F}\) 非迷向，这表明 \(b)\) 蕴含 \(a)\)。

定义 2. --- 设 Q 是 E 上的二次型。若 Q(x) = 0，则称 E 的元素 x 是奇异的（关于 Q 的）。称 E 的子模 F 是：

\begin{enumerate}
\def\labelenumi{\arabic{enumi})}
\item
  奇异的，如果存在 F 中奇异且与 F 正交的元素 \(x \neq 0\)；
\item
  全奇异的，如果 Q 在 F 上的限制为零。
\end{enumerate}

二次模 (E, Q) 的核（§ 3，第 4 小节）由 E° 的奇异元素组成；要使子模 F 是奇异的，必须且只需其核 ≠ \{0\}。由于 Φ(x, y) = Q(x + y) − Q(x) − Q(y)，每个 ≠ \{0\} 的全奇异子模都是奇异的。由于 Φ(x, x) = 2Q(x)，每个奇异向量都是迷向的，且每个奇异（相应地，全奇异）子模都是迷向的（相应地，全迷向的）；当 A 是特征 ≠ 2 的体时，逆命题成立。包含在全奇异子模中的每个子模本身都是全奇异的。两两正交的全奇异子模族之和是全奇异子模。E 的全奇异子模的集合按包含关系排序是归纳的；因此 E 的每个全奇异子模都包含在某个极大全奇异子模之中。

